\section{Composition in one protected task}\label{sec:composition}
We distinguish a certified comparison of experiments from an attained error mechanism. This avoids treating a permitted approximation tolerance as an unavoidable risk lower bound.

\begin{proposition}[Typed causal comparison]\label{prop:composition}
Suppose a parameter-independent causal simulator from a physical experiment $\mathfrak E$ to a target $\mathfrak F$ includes a report transformation, a strategy lift, action compatibility, a stopping handshake and the protected audit readout. Suppose it uses at most $s$ persistent labels and, for every lifted strategy, (i) its complete stopped-and-scored law differs from the target by at most $\varepsilon$ in total variation, and (ii) the expected physical cost is at most $a+b\E\tau$, where $a,b\ge0$ are fixed. Then any target $M$-state observer of expected budget $N$ gives a physical observer with at most $sM$ states, expected cost at most $a+bN$, and risk at most
\[
 R_{\mathfrak F}+L\varepsilon
\]
for a common loss in $[0,L]$. In the same execution law, a terminal decoder perturbation bounded in $L^2$ by $e$ satisfies
\[
 \sqrt{R_{\rm perturbed}}\le\sqrt{R_{\rm exact}}+e.
\]
Independent stage simulation defects may be summed when they are certified conditionally on every surviving history and their expected accumulated sum is bounded. Persistent simulator states multiply; separately stored program bits and workspaces add under a sequential implementation. A reverse risk bound requires a separately certified reverse morphism.
\end{proposition}
\begin{proof}
Run the target label and simulator label as their Cartesian product. The simulator's strategy lift uses no true model parameter, and its terminal handshake precedes the protected audit. The assumed complete-law comparison bounds the expectation of a loss in $[0,L]$ by $L\varepsilon$. The stated cost bound is a separate hypothesis and is not inferred from total variation. Conditional maximal couplings, or telescoping conditional integral differences, bound accumulated bounded-task error by the expected sum of the stage defects. Minkowski's inequality proves the decoder estimate on the same law. A one-way construction says nothing about procedures that do not factor through it, so it cannot prove a reverse minimax comparison.
\end{proof}

Numerical changes to an acceptance rule can change its stopping time. A decoder perturbation alone does not. In particular, closeness of a bounded scored law does not control the mean of an unbounded stopping time. Appendix~\ref{app:compute} gives a separate hazard-margin implementation certificate.

\begin{theorem}[An attained joint resource--error law]\label{thm:joint}
Take either raw model of Section~\ref{sec:realizations}, with $d,\alpha>0$, and augment its protected audit by two independent coordinates. A compulsory preparation call creates a fair sign $U$; a subsequent paid early report is $U$ with probability $1-2\varepsilon$ and the erasure symbol otherwise, where $0\le\varepsilon\le1/2$. A fixed calibration alternative $\zeta\in\{-\delta,\delta\}$ is unobservable in every report. The core experiment and its model are independent of $(U,\zeta)$; the terminal audit is $(Z,U,\zeta)$. The prediction is committed before this audit. Squared losses of the three coordinates are added.

Count preparation, early report and final audit as three physical calls, and put $N=n+3$ with $n\ge1$. Then, for every $M\ge7$ and bounded $\delta$,
\begin{equation}\label{eq:joint-law}
 \mathcal R_{\rm aut}(N,M,\delta,\varepsilon)
 \asymp F_{d,\alpha}(n,M-1)+\delta^2+\varepsilon.
\end{equation}
The constants are independent of $n,M,\delta,\varepsilon$ on the stated range. The early erased-bit experiment has causal deficiency exactly $\varepsilon$ relative to the perfect early-bit interface, whose virtual report must be completed before any core action. Thus the last term is attained, not merely an allowed error ceiling.
\end{theorem}
\begin{proof}
For the lower bound, permit the core controller the entire history, including the early outcome and every fresh coin. The calibration certificate still holds after each of these histories, so the actual core occupation proof with its total expected budget $n$ remains valid. At most $M-1$ terminal core predictions can be used. This gives $cF_{d,\alpha}(n,M-1)$, without assuming that the three early-report lanes receive equal conditional budgets. On the erasure event $U$ remains independent of all subsequently observed core data and routing coins; its Bayes squared loss is one there. Its contribution is at least $2\varepsilon$. Averaging the two calibration alternatives gives at least $\delta^2$, since all report laws are identical under them. These three conditional lower bounds add for the same execution, and averaging is a valid lower bound to the worst calibration alternative and model.

For the upper, put $m=\lfloor(M-4)/3\rfloor\ge1$. Use one initial early-report label, three copies of the core live label, and $3m$ terminal labels. The total is $4+3m\le M$. Preparation occurs before this initial observer label and is nevertheless charged physically; the audit occurs only after terminal commitment and is also charged. Each lane uses the same core program of Theorem~\ref{thm:geometry}. Read out the known bit on the two revealed lanes and zero on the erasure lane; always read out zero for calibration. The risk is at most $CF_{d,\alpha}(n,m)+2\varepsilon+\delta^2$. Since $m$ and $M-1$ differ by a uniformly bounded factor for $M\ge7$, this is \eqref{eq:joint-law}.

For deficiency, regard the two possible bit preparations $U=\pm1$ as the experiment parameter. A causal simulator copies an available bit and uses a fresh fair sign on erasure. Its worst parameter total variation from the perfect early bit is $\varepsilon$. Conversely, under the fair prior, no transformation of an erased bit can guess its sign with error less than $\varepsilon$. The indicator test for an incorrect virtual early report is bounded by total variation, proving the opposite inequality. The prescribed early completion prevents using the final audit to manufacture that report retrospectively. Later core execution can be coupled on the correct virtual-bit event. This proves exact deficiency for the declared complete causal interface.
\end{proof}

The same construction attaches the real $\delta^2$ and $\varepsilon$ terms to the Gaussian and revealing-mass bounds, with the same lane accounting; the term for the core is then the corresponding all-history lower and explicit three-state upper rather than an asserted exact checkpoint value. An additive numerical error requires the independent execution certificate just stated. Theorems~\ref{thm:central}, \ref{thm:geometry}, \ref{thm:reveal} and \ref{thm:joint} therefore distinguish geometric distortion, acquisition uncertainty, realized calibration ambiguity, numerical implementation and causal comparison.
